% ── Rozwiązanie: Exact Subset Sum -- instancja 1 ─────────────────
\subsection{Exact Subset Sum -- instancja 1}

Dane: $S = \{1, 4, 5, 11, 12\}$ (w~kolejności $x_1, \ldots, x_5$) oraz $t = 21$.
Algorytm \textsc{ExactSubsetSum} startuje z~$L_0 = (0)$ i~w~$i$-tej iteracji
scala $L_{i-1}$ z~$L_{i-1} + x_i$ (procedura \textsc{MergeLists} -- posortowane
złączenie bez duplikatów), a~następnie usuwa z~wyniku wszystkie liczby $> t$.
Lista $L_i$ to posortowane sumy podzbiorów $\{x_1, \ldots, x_i\}$ nieprzekraczające
$t$. Element wprowadzony przez przesunięcie $L_{i-1} + x_i$ wyróżniono na~%
\textcolor{accentB}{pomarańczowo}.

\begin{dygresja}
	Przebieg iteracji ($L_{i-1} + x_i$ to składnik przesunięty):
	\begin{enumerate}[nosep]
		\item $x_1 = 1$:\quad
		      $(0) \cup (\textcolor{accentB}{1}) = (0, \textcolor{accentB}{1})$
		      \hfill $L_1 = (0, 1)$
		\item $x_2 = 4$:\quad
		      $(0, 1) \cup (\textcolor{accentB}{4}, \textcolor{accentB}{5})
		      = (0, 1, \textcolor{accentB}{4}, \textcolor{accentB}{5})$
		      \hfill $L_2 = (0, 1, 4, 5)$
		\item $x_3 = 5$:\quad
		      $(0, 1, 4, 5) \cup (\textcolor{accentB}{6}, \textcolor{accentB}{9},
		      \textcolor{accentB}{10})$ \\
		      (przesunięte $5, 6, 9, 10$; $5$ już było)
		      \hfill $L_3 = (0, 1, 4, 5, 6, 9, 10)$
		\item $x_4 = 11$:\quad scalenie z~$(11, 12, 15, 16, 17, 20, 21)$, nic $> 21$
		      \hfill\\
		      $L_4 = (0, 1, 4, 5, 6, 9, 10, \textcolor{accentB}{11},
		      \textcolor{accentB}{12}, \textcolor{accentB}{15},
		      \textcolor{accentB}{16}, \textcolor{accentB}{17},
		      \textcolor{accentB}{20}, \textcolor{accentB}{21})$
		\item $x_5 = 12$:\quad przesunięcie $L_4 + 12$ daje
		      $12, 13, 16, 17, 18, 21, 22, 23, 24, 27, 28, 29, 32, 33$;
		      po~usunięciu $> 21$ dochodzą tylko $\textcolor{accentB}{13}$
		      i~$\textcolor{accentB}{18}$ \\
		      $L_5 = (0, 1, 4, 5, 6, 9, 10, 11, 12, \textcolor{accentB}{13}, 15,
		      16, 17, \textcolor{accentB}{18}, 20, 21)$
	\end{enumerate}
\end{dygresja}

\paragraph{Wynik.}
Największy element listy $L_5$ nieprzekraczający $t = 21$ to
\[
	\boxed{z = 21}
\]
czyli dokładnie $t$ -- istnieje podzbiór o~sumie równej progowi, np.\
$\{4, 5, 12\}$ lub $\{1, 4, 5, 11\}$. Wersja decyzyjna ma więc odpowiedź
\emph{tak}, a~optymalna suma $\leq t$ wynosi $21$.
